\section{Diskussion der Resultate}
\label{sec:resultate}

Die Optimierung hat ihr Ziel erreicht: Das Netz kommt mit deutlich weniger Visuren aus und erfüllt die Anforderungen
an Genauigkeit und Zuverlässigkeit weiterhin. Die eliminierten Visuren stammen fast ausschliesslich aus dem dicht
beobachteten Netzkern im Gebiet Sperre und von wenigen stark überbestimmten Stationen. Die Redundanz wurde damit
dort abgebaut, wo sie keinen nennenswerten Beitrag zur Kontrolle des Netzes mehr leistete.
\bigskip

Die Einsparung im Feld fällt jedoch geringer aus, als die Anzahl der eliminierten Visuren vermuten lässt. Die
Optimierung entfernt einzelne Visuren, aber keine ganzen Stationen. Alle Stationen müssen weiterhin aufgestellt
werden, auch die Station 7020, auf der nur noch drei Visuren verbleiben, und die zweite Aufstellung auf der
Station 2000. Der Zeitgewinn beschränkt sich somit auf das Messen der gestrichenen Visuren und das Aufstellen der
entsprechenden Zielpunkte. Eine grössere Vereinfachung wäre erst möglich, wenn auch ganze Aufstellungen wegfallen
könnten. Dies wurde im Rahmen dieser Optimierung nicht untersucht.
\bigskip

Die Optimierung macht zudem sichtbar, dass das Netz sehr ungleichmässig überbestimmt ist. Einem stark
kontrollierten Netzkern stehen Netzteile gegenüber, die nur schwach oder gar nicht kontrolliert sind. Dazu gehören
die Punkte im Gebiet Brunni und die Punkte auf der Sperre, welche nur von einer Station aus angezielt werden. Die
Optimierung hat an diesen Schwachstellen nichts verändert, da dort keine Visuren entfernt werden konnten. Auch die
Überschreitung der Lageanforderung bei den Punkten TP510 und 6014 lässt sich durch eine Optimierung nicht beheben.
TP510 liegt am Rand des Netzes im Gebiet Tal, 6014 im schwach angebundenen Gebiet Brunni. Für diese Punkte wären
zusätzliche Beobachtungen nötig und nicht weniger. Für die nächste Messkampagne bietet es sich deshalb an, den
im Netzkern eingesparten Aufwand in diese Netzteile zu investieren. Damit würde das Netz insgesamt homogener.
\bigskip

Bei der Interpretation der Resultate sind einige Einschränkungen zu beachten. Die Optimierung wurde an bereits
vorhandenen Messungen durchgeführt. Ob ein Netz, das von Beginn an mit weniger Visuren gemessen wird, dieselben
Resultate liefert, kann erst mit einer neuen Messkampagne überprüft werden. Die Auswahl der Visuren erfolgte zudem
schrittweise und nach eigenem Ermessen. Eine andere Reihenfolge oder Kombination hätte zu einem anderen, möglicherweise
ebenso guten Netz führen können. Schliesslich beruht die Beurteilung auf dem gewählten stochastischen Modell. Würden
die a-priori-Genauigkeiten an die empirischen Werte angepasst, würden sich auch die lokalen Zuverlässigkeiten und
die Koordinatengenauigkeiten verändern, und damit möglicherweise auch die Auswahl der entbehrlichen Visuren.
\bigskip

Für zukünftige Messkampagnen wird empfohlen, die in der Auswertung aufgeführten Visuren wegzulassen und die
Verbindungen der Exzentrumsstation EX7016.4 sowie die Visuren in den Randbereichen unverändert beizubehalten.
Nach der nächsten Messkampagne sollte das optimierte Netz erneut ausgewertet werden, um die hier gewonnenen
Erkenntnisse an neuen Messungen zu bestätigen.
